% Part C. AI-produced by GPT-6 Astra, OpenAI, ChatGPT Pro mode.
\section{Cote 68: exponentials, module functors, and convolution}
\label{sec:C}

\subsection{Conventions and the algebraic exponential}
Throughout this subsection, $K$ is a commutative unital $\mathbb Q$-algebra and $A$ is a commutative unital $K$-algebra. No Noetherian or finite-generation hypothesis is imposed. A \emph{nilpotent} endomorphism has one exponent that annihilates the whole module; a \emph{locally nilpotent} endomorphism has an exponent depending on the vector. This distinction changes the statement on cote~68, scan~039. Statements below use the transcribed mathematical formulas; the doubtful words naming the two rearrangement properties play no role in their definitions.

\begin{lemma}[Proved: finite and elementwise exponential--logarithm]
\label{C:explog}
For an associative unital $\mathbb Q$-algebra $R$, exponential and logarithm give inverse bijections between nilpotent elements $x\in R$ and unipotent elements $1+y\in R$, where $y$ is nilpotent. For a $K$-module $M$ they also give inverse bijections between locally nilpotent $u\in\operatorname{End}_K(M)$ and automorphisms $U\in\operatorname{Aut}_K(M)$ for which $U-1$ is locally nilpotent. The formulas are interpreted vector by vector in the latter case. They commute with every intertwining $K$-linear map.
\end{lemma}
\begin{proof}
In $\mathbb Q[[T]]$, formal differentiation and the constant coefficients prove
$\log(\exp T)=T$ and $\exp(\log(1+T))=1+T$. Modulo $T^n$ these are identities of polynomials, so substitution of a nilpotent element is legitimate even when $R$ is noncommutative: only the commutative subalgebra generated by that one element is used. Moreover $\exp x-1=xP(x)$ with $P(0)=1$, so the two elements have the same nilpotency index. For a locally nilpotent operator, each vector lies in a finite cyclic module $K[u]m$ on which $u$ is nilpotent. The polynomial identities there show that the elementwise formulas are mutually inverse, linear, and compatible on overlapping cyclic modules. The inverse of $\exp u$ is $\exp(-u)$. If $fu=vf$, induction gives $fu^j=v^jf$, and therefore $f\exp u=\exp v f$; the identical argument applies to logarithms of locally unipotent operators.
\end{proof}

\begin{theorem}[Proved: logarithms of unipotent algebra automorphisms]
\label{C:unipotent}
Exponential induces a bijection
\[
\{D\in\operatorname{Der}_K(A):D\text{ is nilpotent}\}
\longrightarrow
\{U\in\operatorname{Aut}_{K\text{-alg}}(A):U-1\text{ is nilpotent}\}.
\]
It preserves the nilpotency index. The same statement holds with both occurrences of ``nilpotent'' replaced by ``locally nilpotent.''
\end{theorem}
\begin{proof}
For a derivation, induction proves
$D^n(ab)=\sum_{j=0}^n\binom njD^j(a)D^{n-j}(b)$.
Summation gives $\exp(D)(ab)=\exp(D)(a)\exp(D)(b)$; all sums are finite for each fixed pair $(a,b)$ in the locally nilpotent case. Also $D(1)=0$ and $D(k)=0$ for $k\in K$, so the exponential is a unital $K$-algebra automorphism.

Conversely put $N=U-1$. If $N^r=0$, define the operator polynomial
\[
U^t=\sum_{j=0}^{r-1}\binom tjN^j\in\operatorname{End}_K(A)[t].
\]
At every nonnegative integer $t=n$ this is the ordinary $n$th power of $U$, hence is multiplicative. For fixed $a,b$, the polynomial $U^t(ab)-U^t(a)U^t(b)\in A[t]$ vanishes at all nonnegative integers. A Vandermonde matrix at sufficiently many distinct rational integers is invertible over $\mathbb Q$, and therefore over $A$; all its coefficients vanish. Differentiating at $t=0$ proves the Leibniz rule for
\[
D=\left.\frac{d}{dt}U^t\right|_{t=0}
 =\sum_{j=1}^{r-1}\frac{(-1)^{j-1}}jN^j=\log U.
\]
The same identities prove $D(k)=0$. Lemma~\ref{C:explog} gives nilpotence, the inverse formulas, and the index assertion. If $N$ is only locally nilpotent, fix $a,b,ab$ and choose $r$ annihilating these three vectors. The preceding argument is then an identity of polynomials on these vectors and proves the same Leibniz rule. No global bound is used in this last argument.
\end{proof}

\begin{theorem}[Proved: the corrected graded-identity statement, C.1(a)]
\label{C:filtered}
Let $I\subset A$ satisfy $I^q=0$, with $q\geq1$, and put $I^0=A$. Define
\begin{align*}
\mathfrak g_I&=\{D\in\operatorname{Der}_K(A):D(I^k)\subset I^{k+1}\text{ for all }k\geq0\},\\
G_I&=\{U\in\operatorname{Aut}_{K\text{-alg}}(A):(U-1)(I^k)\subset I^{k+1}\text{ for all }k\geq0\}.
\end{align*}
Then $G_I$ is exactly the group of $K$-algebra automorphisms preserving $I$ and inducing the identity on $\operatorname{gr}_I(A)$, and
\[
\exp:\mathfrak g_I\xrightarrow{\ \sim\ }G_I,
\qquad
\log:G_I\xrightarrow{\ \sim\ }\mathfrak g_I
\]
are inverse bijections. Both defining series terminate before degree $q$. This supplies a valid replacement for the last corollary of cote~68, scan~046, using the deleted condition on scan~045.
\end{theorem}
\begin{proof}
The endomorphisms $T$ satisfying $T(I^k)\subset I^{k+1}$ form a nonunital associative algebra $F^1$ with $(F^1)^q=0$: a product of $j$ such endomorphisms carries $I^k$ into $I^{k+j}$. Every nonconstant polynomial in $T\in F^1$ remains in $F^1$. Consequently $\exp D-1\in F^1$ and $\log U\in F^1$. Theorem~\ref{C:unipotent} identifies these polynomials as an algebra automorphism and a derivation, respectively. Their inverse identities follow from Lemma~\ref{C:explog}. If $U-1\in F^1$, the finite geometric series for $U^{-1}$ shows that $U^{-1}-1\in F^1$ too; thus $U(I)=I$. The condition on each quotient $I^k/I^{k+1}$ is precisely the displayed definition of $G_I$. The case $I=0$ gives the singleton group and the zero derivation.
\end{proof}

\subsection{Differential automorphisms: an exact classification}
We use differential operators in the iterated-commutator sense of Grothendieck and Dieudonn\'e (1967), EGA~IV, Proposition~16.8.8 (printed pp.~42--43), equivalently the recursive definition in the Stacks Project, Definition~10.133.1, tag~09CH. Both sources were opened and the equivalence checked~\cite{C_EGAIV4,C_Stacks09CH}. Write $m_a$ for multiplication by $a\in A$.

\begin{theorem}[Proved: differential order and the difference ideal, C.1(b)]
\label{C:differential}
For $U\in\operatorname{Aut}_{K\text{-alg}}(A)$ put
$J_U=(U(a)-a:a\in A)\subset A$.
For every integer $n\geq0$, $U$ is a differential operator of order at most $n$ if and only if $J_U^{n+1}=0$. In particular the differential automorphisms are exactly the automorphisms inducing the identity modulo some nilpotent ideal. Among them, exactly those with $U-1$ globally nilpotent are exponentials of nilpotent derivations.
\end{theorem}
\begin{proof}
The identity
\[
[U,m_a]=m_{U(a)-a}U
\]
and the commutativity of multiplication operators give, by induction,
\[
[\cdots[[U,m_{a_0}],m_{a_1}],\ldots,m_{a_n}]
 =m_{\prod_{i=0}^n(U(a_i)-a_i)}U.
\]
An iterated commutator vanishes precisely when its multiplying coefficient is zero, as evaluation at $1$ shows. The recursive definition of differential order therefore says exactly that every such product is zero, which is equivalent to $J_U^{n+1}=0$. If $(U-1)(A)\subset I$ and $I$ is nilpotent, then $J_U\subset I$; conversely we may take $I=J_U$. The final assertion is Theorem~\ref{C:unipotent}. For completeness, if $D^r=0$, each $D^j$ has differential order at most $j$, by the Leibniz rule or Stacks Lemma~10.133.2; hence $\exp D$ has order at most $r-1$ and $J_{\exp D}^{r}=0$.
\end{proof}

\begin{theorem}[Proved: the conormal criterion, without finite type]
\label{C:conormal}
Suppose $J=J_U$ and $J^q=0$. The ideal $J$ is $U$-stable, and the induced map $U_1$ on $J/J^2$ is linear over $B=A/J$. The following are equivalent:
\begin{enumerate}
\item $U$ is the exponential of a globally nilpotent $K$-derivation;
\item $U-1$ is globally nilpotent;
\item $U_1-1$ is globally nilpotent on $J/J^2$.
\end{enumerate}
If $(U_1-1)^r=0$, then
\[
(U-1)^{\,q+(r-1)q(q-1)/2}=0.
\]
Thus the other differential automorphisms are exactly those whose conormal action has no uniform unipotence exponent. This includes conormal actions that are locally unipotent with unbounded exponents.
\end{theorem}
\begin{proof}
The image under $U$ of a generator $U(a)-a$ is $U^2(a)-U(a)$, and the same calculation with $U^{-1}$ gives $U(J)=J$. Since $U$ is the identity modulo $J$, its action on $J/J^2$ is $B$-linear. The implications (1)$\Leftrightarrow$(2)$\Rightarrow$(3) have already been proved.

For the converse, multiplication induces a surjective $B$-linear map
$(J/J^2)^{\otimes_B k}\twoheadrightarrow J^k/J^{k+1}$, equivariant for $U$. Write $U_1=1+N$, $N^r=0$. On the $k$th tensor power, the induced automorphism is $\prod_{j=1}^k(1+N_j)$, where the $N_j$ commute and $N_j^r=0$. Every monomial of total degree $k(r-1)+1$ in these operators vanishes, by the pigeonhole principle. Consequently
\[
(U_k-1)^{k(r-1)+1}=0\quad\text{on }J^k/J^{k+1}.
\]
The exponent for the quotient $A/J$ is $1$. Applying these bounds successively along the finite filtration $A\supset J\supset\cdots\supset J^q=0$ yields exponent
$\sum_{k=0}^{q-1}(k(r-1)+1)=q+(r-1)q(q-1)/2$ on $A$. This proof uses a surjection of tensor powers, not finite generation or a choice of basis.
\end{proof}

\begin{proposition}[Proved: all automorphisms of a split square-zero thickening]
\label{C:squarezero}
Let $B$ be a commutative $K$-algebra, $M$ any $B$-module, and $A=B\oplus M$ with multiplication
$(b,m)(c,n)=(bc,bn+cm)$. Every $K$-algebra automorphism of $A$ inducing the identity on $B=A/M$ is uniquely
\[
U_{d,\sigma}(b,m)=(b,d(b)+\sigma(m)),
\quad d\in\operatorname{Der}_K(B,M),\quad\sigma\in\operatorname{Aut}_B(M).
\]
Composition is $(d,\sigma)(e,\tau)=(d+\sigma e,\sigma\tau)$. All these automorphisms have differential order at most one. They are exponentials of globally nilpotent derivations precisely when $\sigma-1$ is globally nilpotent.
\end{proposition}
\begin{proof}
Since the automorphism preserves $M$ and induces the identity on the quotient, its restriction to $M$ is an additive bijection $\sigma$, while its value on $(b,0)$ is $(b,d(b))$. Multiplicativity on $(b,0)(0,m)$ proves $B$-linearity of $\sigma$, and multiplicativity on $(b,0)(c,0)$ proves the Leibniz rule for $d$. The converse and the composition formula follow by direct multiplication. The difference ideal lies in $M$ and $M^2=0$, so Theorem~\ref{C:differential} applies. For $N=U-1$ one computes
$N^k(b,m)=(0,(\sigma-1)^{k-1}d(b)+(\sigma-1)^k m)$ for $k\geq1$. Therefore $N$ is globally nilpotent exactly when $\sigma-1$ is, and Theorem~\ref{C:unipotent} finishes the proof.
\end{proof}

\begin{proposition}[Refuted, with corrected statements: scan~045 and scan~046]
\label{C:counterexamples}
On $A=\mathbb Q[\varepsilon]/(\varepsilon^2)$, the automorphism $U(\varepsilon)=2\varepsilon$ is differential of order one but is not an exponential of a nilpotent derivation; it acts as multiplication by $2$ on $J_U/J_U^2$. On $A'=\mathbb Q[\varepsilon]/(\varepsilon^3)$, the derivation $D=\varepsilon\,d/d\varepsilon$ has $D(A')\subset I=(\varepsilon)$ but is not locally nilpotent. Thus identity on $A/I$ does not imply identity on $\operatorname{gr}_I(A)$, and merely requiring $D(A)\subset I$ does not make the exponential finite. The correct replacements are Theorems~\ref{C:filtered}, \ref{C:differential}, and \ref{C:conormal}.
\end{proposition}
\begin{proof}
In the first example $J_U=(\varepsilon)$ and $J_U^2=0$, while $(U-1)^n(\varepsilon)=\varepsilon$ for every positive $n$. In the second example $D(\varepsilon)=\varepsilon$, so $D^n(\varepsilon)=\varepsilon$ for every positive $n$; also $D(\varepsilon^2)=2\varepsilon^2$, so $D$ descends to $A'$. The asserted ideal inclusions and failures follow. Exact matrix calculations for these examples are attached in the program and output files. The prior study attributes the first scaling example also to Michel Hua's project; we retain that attribution and independently verify it here.
\end{proof}

\subsection{Intertwining and additive group actions}
\begin{theorem}[Proved: C.1(c), including infinite modules]
\label{C:intertwine}
Let $M,N$ be arbitrary $K$-modules, $u\in\operatorname{End}_K(M)$ and $v\in\operatorname{End}_K(N)$ locally nilpotent, and $f\in\operatorname{Hom}_K(M,N)$. Then
\[
f\exp(u)=\exp(v)f\quad\Longleftrightarrow\quad fu=vf.
\]
The equivalence also identifies intertwining at time $1$ with intertwining for every parameter over every commutative $K$-algebra, and with intertwining of the formal additive group actions.
\end{theorem}
\begin{proof}
The reverse implication follows by induction and the elementwise finite exponential series. For the forward implication put $P=\exp u-1$, $Q=\exp v-1$. Then $fP=Qf$, hence $fP^j=Q^jf$ for every $j$. On each vector the logarithm series terminates, so
$f\log(1+P)=\log(1+Q)f$. Lemma~\ref{C:explog} identifies these logarithms with $u$ and $v$. Tensoring the intertwining identity with a $K$-algebra $R$ proves it after base change and then for $\exp(\lambda u_R)$ for every $\lambda\in R$. Conversely the universal parameter over $K[t]/(t^2)$ recovers $fu=vf$ from its coefficient of $t$.
\end{proof}

\begin{remark}[Proved correction of the proof on scan~041]
\label{C:proof-correction}
In the manuscript's alternative proof, $\delta=vf-fu$ lies in $\operatorname{Hom}_K(M,N)$, and the commuting operators are $T(\phi)=v\phi$ and $T'(\phi)=\phi u$ on this Hom module. The written choice $U=\operatorname{End}_K(M)$ is incorrectly typed unless $M=N$. Expanding the exponential difference gives
\[
0=\delta+\sum_{i+j\geq1}\frac{v^i\delta u^j}{(i+j+1)!},
\]
so the coefficients in $\delta=\sum c_{ij}v^i\delta u^j$ are $c_{ij}=-1/(i+j+1)!$. This corrects both the sign and factorial index. The elementwise-logarithm proof above avoids requiring a globally nilpotent operator on the entire Hom module when $u$ and $v$ are only locally nilpotent.
\end{remark}

\begin{theorem}[Proved: the exact module version of scan~039]
\label{C:module-actions}
For any $K$-module $M$, let $\mathcal M(R)=M\otimes_K R$ and
$\widehat{\mathbb G}_a(R)=\operatorname{Nil}(R)$. Natural $R$-linear additive-group actions on $\mathcal M$ satisfy
\begin{align*}
\operatorname{End}_K(M)&\simeq
\operatorname{Hom}_{\mathrm{grp}}(\widehat{\mathbb G}_a,\operatorname{Aut}(\mathcal M)),\\
\operatorname{LNil}(\operatorname{End}_K(M))&\simeq
\operatorname{Hom}_{\mathrm{grp}}(\mathbb G_a,\operatorname{Aut}(\mathcal M)),
\end{align*}
by $u\mapsto[\lambda\mapsto\exp(\lambda u_R)]$.
Here $\operatorname{LNil}$ means locally nilpotent endomorphisms. Restricting the second map to globally nilpotent endomorphisms is bijective exactly when every locally nilpotent endomorphism of $M$ is globally nilpotent. Finite generation of $M$ is sufficient, but is not necessary.
\end{theorem}
\begin{proof}
For a formal action, evaluate at the universal parameter $t$ over $K[t]/(t^{r+1})$ and apply it to $m\otimes1$. The unique expansion in the free $K$-basis $1,t,\ldots,t^r$ gives $K$-linear coefficients $u_j(m)$. Naturality in the truncation maps makes these coefficients independent of $r$. Naturality under $t\mapsto\lambda$ proves that they determine the action on all scalar extensions. The group law, evaluated over $K[s,t]/(s,t)^{r+1}$, gives
\[
u_0=1,\qquad \binom{i+j}{i}u_{i+j}=u_i u_j.
\]
Thus $u_j=u_1^j/j!$. Conversely these coefficients give the formal exponential for every $u_1$, since $\lambda$ is nilpotent.

For an ordinary action the same evaluation over $K[t]$ gives, for each $m$, a polynomial with finitely many coefficients. The same group law implies $u_1^j(m)=0$ for large $j$, precisely local nilpotence. Conversely local nilpotence makes the expression a polynomial on every vector, and therefore defines a natural action. A finite generating set supplies a uniform nilpotence exponent if $M$ is finitely generated.

For nonnecessity, take $K=\mathbb Q[x]$ and $M=\bigoplus_{n\geq0}K/(x-n)$. Distinct summands have zero $K$-linear Hom, so every endomorphism is diagonal by scalars in $\mathbb Q$. A locally nilpotent diagonal endomorphism is zero. This $M$ is not finitely generated, but the stated restricted map is still bijective. In contrast, for $K=\mathbb Q$, $M=\mathbb Q[x]$, differentiation is locally nilpotent but not globally nilpotent, and its action is translation of polynomials. This verifies the needed correction without asserting finite generation to be necessary.
\end{proof}

\subsection{Arbitrary module functors: the full obstruction criterion, C.1(d)}
An $\mathcal O$-module functor $\mathcal M$ assigns an $R$-module to each commutative $K$-algebra $R$ and a semilinear transition map to each algebra homomorphism. Write $E=\operatorname{End}_{\mathcal O}(\mathcal M)$ for its ring of natural $R$-linear endomorphisms. An action here means a natural transformation $\mathbb G_a\times\mathcal M\to\mathcal M$ (or its formal analogue), additive in the group parameter and $R$-linear in the module argument. This explicitly specifies the interpretation of $\operatorname{Aut}_K(\mathcal M)$ in scan~044; it is not an assumption that $\mathcal M$ is a tensor-product functor.

Let $\mathcal N(R)=\operatorname{Nil}(R)$, an ideal and therefore an $R$-module. Tensor products of module functors below are taken objectwise. Define the explicit $K$-linear map
\[
\delta:E\longrightarrow
H=\operatorname{Hom}_{\mathcal O}(\mathcal N\otimes_{\mathcal O}\mathcal M,\mathcal M),
\qquad
\delta(u)_R(\lambda\otimes m)=\lambda u_R(m).
\]

\begin{theorem}[Proved: exact images, fibers, and bijectivity criteria]
\label{C:functor-criterion}
For an arbitrary $\mathcal O$-module functor:
\begin{enumerate}
\item The map $u\mapsto\bar e_u$ from $E$ to formal additive-group actions has fibers exactly the additive cosets of $\ker\delta$. A formal action $\rho$ belongs to its image if and only if all three following tests hold: each $\rho_R(\lambda)$ is unipotent for nilpotent $\lambda$; its finite logarithm is $R$-linear in the parameter $\lambda\in\operatorname{Nil}(R)$; and the resulting morphism $\ell_\rho\in H$ has zero class in $\operatorname{coker}\delta$. In particular the formal exponential map is bijective if and only if $\ker\delta=0$ and these three tests hold for every formal action.
\item The map from globally nilpotent elements of $E$ to ordinary additive-group actions is always injective. Its image consists exactly of those actions for which there is one integer $r$ with $(\rho_R(\lambda)-1)^r=0$ for all $R,\lambda$, and whose logarithms satisfy
\[
\log\rho_R(a\lambda)=a\log\rho_R(\lambda)
\quad(a,\lambda\in R).
\]
It is bijective exactly when every ordinary action passes these two tests.
\end{enumerate}
These criteria distinguish failure of uniqueness, failure of unipotence, failure of scalar linearity, and failure to divide the formal parameter. They apply without representability, exactness, or finiteness assumptions on $\mathcal M$.
\end{theorem}
\begin{proof}
If $\rho=\bar e_u$, then for $\lambda^r=0$ the operator $\lambda u_R$ has $r$th power zero. Lemma~\ref{C:explog} gives $\log\rho_R(\lambda)=\lambda u_R$, which passes the three tests. Conversely, assume the tests and choose $u\in E$ with $\delta(u)=\ell_\rho$. Then $\log\rho_R(\lambda)=\lambda u_R$ for all $R,\lambda$, and exponential--logarithm inversion gives $\rho=\bar e_u$.

To justify that the logarithms used here form an additive natural map, note that the operators $\rho_R(\lambda)$ commute for different parameters, by the group law. Two commuting unipotent operators generate a commutative algebra whose ideal generated by their differences from $1$ is nilpotent. The finite identity $\log(UV)=\log U+\log V$ in that algebra therefore proves additivity. Naturality follows because every transition map intertwines the corresponding operators and their finite logarithms. The second test gives $R$-linearity in $\lambda$, while linearity in $m$ is automatic, so $\ell_\rho$ is exactly an element of $H$. Equality $\bar e_u=\bar e_v$ is equivalent, after finite logarithms, to $\delta(u-v)=0$, proving the fiber assertion.

For an ordinary action in the image, $u^r=0$ supplies the uniform exponent, and its logarithm is $\lambda u_R$, so the tests are necessary. If they hold, put $u_R=\log\rho_R(1)$. This defines a natural endomorphism $u\in E$, and $u^r=0$ because logarithm is a polynomial without constant term in $\rho(1)-1$. Scalar linearity with parameter $1$ gives $\log\rho_R(\lambda)=\lambda u_R$, proving sufficiency. Finally evaluation at $1$ followed by the logarithm recovers $u$, proving injectivity.
\end{proof}

\begin{proposition}[Refuted: universal bijectivity in the question on scan~044]
\label{C:functor-counterexamples}
Both maps asked about on cote~68, scan~044, fail to be bijective for general module functors. The manuscript poses a question here, not an assertion of bijectivity.
\begin{enumerate}
\item For $K=\mathbb Q$ and $\mathcal M(R)=R_{\mathrm{red}}$, the distinct endomorphisms $0,1\in E$ have the same formal exponential.
\item For $K=\mathbb Q$ and $\mathcal M(R)=R\oplus\Omega_{R/\mathbb Q}$, the formula
\[
\rho_R(\lambda)(a,\omega)=(a,\omega+a\,d\lambda)
\]
defines an ordinary additive-group action, with $(\rho_R(\lambda)-1)^2=0$ for every $R,\lambda$. Neither it nor its restriction to the formal additive group is an exponential of the required form.
\end{enumerate}
\end{proposition}
\begin{proof}
For (1), every nilpotent scalar acts by zero on $R_{\mathrm{red}}$, so $\exp(\lambda u_R)=1$ for every endomorphism $u$ and nilpotent $\lambda$. The endomorphisms $0$ and $1$ differ already at $R=\mathbb Q$.

For (2), the naturality of K\"ahler differentials makes $\mathcal M$ an $\mathcal O$-module functor and makes the displayed maps natural. Each is $R$-linear in $(a,\omega)$; additivity of $d$ proves the action law, and the inverse uses $-\lambda$. The difference from the identity has square zero. If the ordinary action were $e_u$ with globally nilpotent $u$, evaluation at $\lambda=1$ would give $\exp u=1$, hence $u=0$. But at $R=\mathbb Q[t]$, $\rho_R(t)(1,0)=(1,dt)\ne(1,0)$.

For the formal restriction, use $R=\mathbb Q[\varepsilon]/(\varepsilon^2)$. Here $\Omega_{R/\mathbb Q}=\mathbb Q\,d\varepsilon$, with $\varepsilon\,d\varepsilon=0$ and $d\varepsilon\ne0$. If the action were $\bar e_u$, then
$\rho_R(\varepsilon)-1=\varepsilon u_R$, since $\varepsilon^2=0$. The second component of $\varepsilon u_R(1,0)$ is zero, whereas that of $(\rho_R(\varepsilon)-1)(1,0)$ is $d\varepsilon$. This contradiction proves the claim. The obstruction is the failure of the logarithmic parameter map to lie in the image of $\delta$; its scalar linearity also fails, for example $d(t\varepsilon)\ne t\,d\varepsilon$ over $\mathbb Q[t,\varepsilon]/(\varepsilon^2)$.
\end{proof}

\begin{theorem}[Proved: a structural base-extension criterion, including internal Hom]
\label{C:functor-basechange}
Let $\mathcal M$ be an arbitrary $\mathcal O$-module functor over the commutative $\mathbb Q$-algebra $K$.
\begin{enumerate}
\item Suppose that for every $K$-algebra $R$ and every $n\geq1$, the canonical map
\[
 \mathcal M(R)\otimes_R R[t]/(t^{n+1})
       \longrightarrow\mathcal M(R[t]/(t^{n+1}))
\]
is an isomorphism. Then $u\mapsto\bar e_u$ is bijective.
\item Suppose instead that for every $R$, the canonical map
\[
 \mathcal M(R)\otimes_R R[t]\longrightarrow\mathcal M(R[t])
\]
is an isomorphism. Then ordinary additive-group actions correspond exactly to natural endomorphisms $u$ for which each $u_R$ is elementwise locally nilpotent. The map from the globally nilpotent elements of $E$ is therefore bijective exactly when all these natural locally nilpotent endomorphisms are globally nilpotent in $E$.
\end{enumerate}
The first hypothesis holds for
\[
 \mathcal M(R)=\operatorname{Hom}_R(P\otimes_KR,Q\otimes_KR)
\]
with arbitrary $K$-modules $P,Q$, without finite generation or projectivity, and for arbitrary product functors $\mathcal M(R)=\prod_{j\in J}R$. Both hypotheses hold for every tensor-product functor $M\otimes_KR$. For the internal Hom functor the second hypothesis also holds when $P$ is finitely generated.
\end{theorem}
\begin{proof}
Under the first hypothesis, evaluate a formal action at the universal parameter $t$ over $R[t]/(t^{n+1})$, on an element $m\in\mathcal M(R)$. The assumed isomorphism gives a unique expression
\[
 \rho_{R[t]/(t^{n+1})}(t)(m\otimes1)
       =\sum_{j=0}^n u_{j,R}(m)t^j.
\]
Linearity of the action gives $R$-linear coefficients, naturality in $R$ makes them natural endomorphisms, and the quotient maps between truncated polynomial algebras make them independent of $n$. The group law can be compared coefficient by coefficient in
$R[s,t]/(s^{n+1},t^{n+1})$: the required direct-sum description of $\mathcal M$ on that ring follows by applying the stated hypothesis twice. Naturality from $R[z]/(z^{2n+1})$ under $z\mapsto s+t$ makes the left side of the group law meaningful there and yields
\[
 u_{0,R}=1,\qquad \binom{i+j}{i}u_{i+j,R}=u_{i,R}u_{j,R}
 \quad(0\le i,j\le n).
\]
Thus $u_{j,R}=u_{1,R}^j/j!$. Write $u=u_1\in E$. For any nilpotent $\lambda\in R$, choose $n$ with $\lambda^{n+1}=0$ and apply naturality to the $R$-algebra map $R[t]/(t^{n+1})\to R$, $t\mapsto\lambda$. It gives $\rho_R(\lambda)=\exp(\lambda u_R)$. The coefficient of $t$ over $R[t]/(t^2)$ uniquely recovers $u_R$ for every $R$, proving both surjectivity and injectivity.

Under the second hypothesis, use $R[t]$ instead. The image of each $m\otimes1$ has only finitely many nonzero polynomial coefficients. The same coefficient comparison over $R[s,t]$ gives $u_j=u_1^j/j!$, hence $u_R^j(m)=0$ for all large $j$, depending on $R,m$. Conversely such pointwise local nilpotence makes every exponential series finite on its argument, gives natural $R$-linear maps, and verifies the action law and inverse by finite polynomial identities. The assertion about the globally nilpotent subdomain follows.

Finally, if $S=R[t]/(t^{n+1})$, then $S$ is finite free over $R$, and adjunction gives
\begin{align*}
 \operatorname{Hom}_S(P_R\otimes_RS,Q_R\otimes_RS)
 &\simeq\operatorname{Hom}_R(P_R,Q_R\otimes_RS)\\
 &\simeq\operatorname{Hom}_R(P_R,Q_R)\otimes_RS.
\end{align*}
The second equality only uses a finite direct sum, which Hom commutes with for every $P_R$; it assumes no finite presentation. Finite direct sums also commute with arbitrary products, proving the product-functor assertion. Tensor-product functors satisfy both comparisons by associativity of tensor product. If $P$ is finitely generated, each $P_R$ is finitely generated and Hom from it commutes also with the direct sum $Q_R\otimes_RR[t]=\bigoplus_{j\ge0}Q_Rt^j$: the images of finitely many generators occupy finitely many summands. This proves the last statement.
\end{proof}

\begin{proposition}[Proved: the formal unipotence and parameter-division obstructions occur]
\label{C:functor-separate-obstructions}
The distinct tests in Theorem~\ref{C:functor-criterion} cannot in general be omitted.
\begin{enumerate}
\item There is a natural formal additive-group action whose value at a square-zero scalar is not even locally unipotent.
\item There is a natural formal action which is uniformly unipotent of exponent two and whose logarithm is $R$-linear in its nilpotent scalar parameter, but which is not a formal exponential because its logarithm has nonzero class in $\operatorname{coker}\delta$.
\end{enumerate}
\end{proposition}
\begin{proof}
For (1), take $K=\mathbb Q$ and the free $R$-module
\[
 \mathcal M(R)=\bigoplus_{\mu\in\operatorname{Nil}(R)}R[\mu].
\]
A map $\alpha:R\to S$ sends $a[\mu]$ to $\alpha(a)[\alpha(\mu)]$; this specifies its functorial semilinear transition. Translation of the basis,
$\rho_R(\lambda)[\mu]=[\mu+\lambda]$, is a natural $R$-linear formal additive-group action. At $R=\mathbb Q[\varepsilon]/(\varepsilon^2)$, the elements $j\varepsilon$ are pairwise distinct for distinct integers $j$. Therefore
\[
 (\rho_R(\varepsilon)-1)^m[0]
   =\sum_{j=0}^m(-1)^{m-j}\binom mj[j\varepsilon]\ne0
\]
for every $m$. In contrast, the formal exponential of any endomorphism at $\varepsilon$ is $1+\varepsilon u_R$ and has square-zero difference from the identity.

For (2), take $\mathcal M(R)=R\oplus\operatorname{Nil}(R)$, with its evident transitions, and define
\[
 \rho_R(\lambda)(a,n)=(a,n+a\lambda),\qquad
 \lambda\in\operatorname{Nil}(R).
\]
This is natural and $R$-linear in $(a,n)$. Its difference from the identity has square zero, its products add the parameter, and its logarithm is exactly $(a,n)\mapsto(0,a\lambda)$, which is $R$-linear in $\lambda$. Thus it defines an element of $H$ and passes the first two formal-image tests. But over the preceding dual-number ring, every element of $\operatorname{Nil}(R)=\mathbb Q\varepsilon$ is annihilated by $\varepsilon$. For any $R$-linear endomorphism $u_R$ of $R\oplus\operatorname{Nil}(R)$, the second component of $\varepsilon u_R(1,0)$ is consequently zero. The second component of $\log\rho_R(\varepsilon)(1,0)$ is $\varepsilon\ne0$. Hence this logarithm is not in $\delta(E)$. This proves the nonzero cokernel obstruction even after unipotence and scalar linearity have both been verified. The attached \texttt{code/C\_functor\_obstructions.py} records the exact three-dimensional dual-number matrices and the linear functional separating this logarithm from all products $\varepsilon u_R$.
\end{proof}

\paragraph{What is determined for arbitrary functors.}
Theorem~\ref{C:functor-criterion} gives necessary and sufficient conditions for each action to lie in the image, all possible preimages, and hence exact criteria for either map to be bijective. Theorem~\ref{C:functor-basechange} gives conditions stated directly on the functor under which those criteria simplify, including the internal Hom functors that motivate scan~043. The counterexamples separately show loss of uniqueness, nonunipotent actions, failure of scalar linearity, and failure to divide a nilpotent parameter. No claim is made that the sufficient base-extension properties are necessary, or that arbitrary module functors have been classified up to isomorphism.

\begin{proposition}[Proved: globalization and the meaning of local nilpotence]
\label{C:sheaf-globalization}
Let $X$ be a scheme over $\mathbb Q$, and let $D$ be a derivation of the sheaf $\mathcal O_X$. Suppose each point has an open neighborhood $V$ and an integer $r$ such that $D^r=0$ as an endomorphism of $\mathcal O_X|_V$. The local finite exponentials glue to an automorphism $E_D$ of $X$ inducing the identity on $X_{\mathrm{red}}$. If $X\to S$ is a morphism and $D$ is an $S$-derivation, this is an $S$-automorphism. If $D,D'$ satisfy this hypothesis and commute, then $D+D'$ satisfies it and
\[
 E_0=\operatorname{id}_X,\qquad E_{D+D'}=E_D\circ E_{D'}.
\]
Thus the plus sign between automorphisms in the transcription of scan~044 must be composition. Elementwise local nilpotence on the ring of one affine open is a weaker hypothesis and does not imply identity on its reduction.
\end{proposition}
\begin{proof}
On each such neighborhood, the exponential is an automorphism of sheaves of $\mathbb Q$-algebras by Theorem~\ref{C:unipotent}, with inverse $\exp(-D)$. The formulas agree on overlaps, because they are the same finite series with zero terms appended. For any smaller affine open $W\subset V$, the induced endomorphism $D_W$ of $A=\mathcal O_X(W)$ has $D_W^r=0$. Theorem~\ref{C:differential} gives $J_{\exp D_W}^r=0$, so the exponential induces the identity on $A_{\mathrm{red}}$. In particular it preserves every prime ideal and its maps on local rings are local. It therefore defines a scheme automorphism on the identity underlying topological map and acts identically on the reduction. If $D$ vanishes on the image of $\mathcal O_S$, every nonconstant exponential term vanishes there as well.

Locally choose bounds $D^r=0$, $(D')^s=0$. Commutativity and the binomial formula give $(D+D')^{r+s-1}=0$, while multiplication of the finite exponential series gives $\exp(D+D')=\exp D\,\exp D'$. Pullbacks reverse compositions, but these two pullbacks commute, giving the displayed scheme identity. Finally $d/dx$ on $\mathbb Q[x]$ is elementwise locally nilpotent and exponentiates to $f(x)\mapsto f(x+1)$, a nontrivial automorphism of the reduced affine line. It is not nilpotent with a uniform bound, nor does its exponential define a sheaf automorphism on the identity topological map. This supplies the precise distinction required in scan~044.
\end{proof}

\subsection{Convolution and the exact content of the system axioms}
Let $G$ be any discrete group, with identity $e$. Let $\mathbb R_+[G]$ denote the nonnegative finitely supported real functions on $G$, and put
\[
(f*g)(x)=\sum_{y\in G}f(y)g(y^{-1}x),\qquad
\check f(x)=f(x^{-1}),\qquad L_s f(x)=f(s^{-1}x).
\]
For $k\geq1$ not exceeding $|G|$ in the finite case, let $p_k(f)$ be the sum of the $k$ largest values, padding by zeros. We write $f'\preceq f$ when $p_k(f')\leq p_k(f)$ for every $k$, and $f'\sim f$ when their multisets of nonzero values agree. Set
$R(f,g)=\sum_{j\geq1}f_j^\downarrow g_j^\downarrow$.
Following the image-checked terminology of scans~020--021, a pair is \emph{pr\'etass\'e} if $\|f*g\|_\infty=R(f,g)$; it is \emph{tass\'e} if $f'\preceq f$, $g'\preceq g$ always imply $f'*g'\preceq f*g$.

\begin{lemma}[Proved: rearrangement, partial sums, and convex reduction]
\label{C:majorization}
For nonnegative finitely supported functions on a set, the maximum scalar product over rearrangements is $R(f,g)$. It is attained exactly when all positive level sets of the two functions are mutually nested. Moreover $f'\preceq f$ is equivalent to
$R(f',h)\leq R(f,h)$ for every nonnegative finitely supported $h$, and to this inequality only for indicator functions of finite subsets. On a finite universe the set $\{f':f'\preceq f\}$ is the convex hull of the functions lying coordinatewise below a permutation of $f$.
\end{lemma}
\begin{proof}
Write a nonnegative function as the finite layer-cake sum
$f=\sum_i\alpha_i1_{A_i}$, where $\alpha_i>0$ and the nonempty $A_i$ form its nested positive level sets; likewise $g=\sum_j\beta_j1_{B_j}$. Since $|A_i\cap B_j|\leq\min(|A_i|,|B_j|)$, summation proves the rearrangement bound. All terms have positive coefficients, so equality holds exactly when every pair of level sets is comparable by inclusion. Aligning the decreasing orders realizes equality. Taking $h=1_H$ gives $R(f,h)=p_{|H|}(f)$; summation by parts against the decreasing values of arbitrary $h$ proves the converse.

For the last assertion, let $Q_f$ be the stated convex hull. The partial sums $p_k$ are convex and coordinatewise increasing, so $Q_f$ lies in the majorization set. Conversely, for any real coefficient vector $w$,
\[
\max_{z\in Q_f}w\cdot z
 =\sum_i(w^+)^\downarrow_i f_i^\downarrow.
\]
If $x\preceq f$, rearrangement and summation by parts give
$w\cdot x\leq\sum_i(w^+)^\downarrow_i x_i^\downarrow
\leq\sum_i(w^+)^\downarrow_i f_i^\downarrow$.
The finite-dimensional separating-hyperplane theorem now gives $x\in Q_f$. For applications on an infinite group, all functions under discussion have finite support; a finite universe containing these supports and enough zero coordinates gives the same reduction. In particular a positive bilinear expression that is separately convex after applying $p_k$ need only be maximized on rearrangements, since coordinatewise reductions cannot increase it.
\end{proof}

We inspected scan~017 directly, including its upside-down lower half. Its displayed data are a homomorphism $\rho:\widetilde G\to G$, an identification $i:G[2]\simeq\ker\rho$, and chains $\mathcal A(\tilde g)$ of finite subsets of $G$. Here $G[2]=\{x\in G:x^2=e\}$ when this set is a subgroup. Put
\[
C(\tilde g)=\left\{\sum_{A\in F}c_A1_A:
 F\subset\mathcal A(\tilde g)\text{ finite},\ c_A\geq0\right\},
\quad S=\bigcup_{\tilde g}C(\tilde g).
\]
The axioms read: each $\mathcal A(\tilde g)$ is a chain; $A=\rho(\tilde g)A^{-1}$ for its members; convolution of two generating indicators belongs to $C(\tilde g\tilde h)$; translation by $u\in G[2]$ takes $\mathcal A(\tilde g)$ to $\mathcal A(i(u)\tilde g)$; and two chains become jointly nested after a left translation of one. The last two requirements are both labelled d) in the manuscript. In c), the second family is written with a letter transcribed as $\mathcal B$ although only $\mathcal A$ has been defined. We use the natural same-family reading. Our counterexample also works if two families were intended, by taking them identical. No conclusion below depends on deciding between those two readings.

The \emph{cone-wise span} $S$ is convolution-stable by distributivity. If instead ``the family they span'' means the single positive cone generated by all families together, the counterexample below still belongs to that larger cone and still refutes Condition~K. The positive results concerning pairs in $S$ use the cone-wise definition just displayed.

\begin{theorem}[Proved: what a)--d) actually imply]
\label{C:pretasse-cones}
For the preceding system, every pair $(f,g)\in S^2$ is pr\'etass\'e. In particular the two-factor peak inequality
\[
f'\preceq f,\quad g'\preceq g
\quad\Longrightarrow\quad
\|f'*g'\|_\infty\leq\|f*g\|_\infty
\]
holds. This does not imply the corresponding inequality for three factors, or two-factor majorization.
\end{theorem}
\begin{proof}
Choose labels with $f\in C(\tilde a)$, $g\in C(\tilde b)$, and set $b=\rho(\tilde b)$. The reflection condition gives $\check g=L_{b^{-1}}g$. Joint nesting gives $s\in G$ such that every level set of $f$ and of $L_sg$ is nested. Thus Lemma~\ref{C:majorization} gives
\[
(f*g)(sb)=\langle f,L_{sb}\check g\rangle
 =\langle f,L_sg\rangle=R(f,g).
\]
Every value of every convolution of rearrangements is at most this rearrangement bound. Weak majorization only decreases the bound, by the same lemma. The final two nonimplications are proved by the next counterexample.
\end{proof}

\begin{theorem}[Refuted: C.2(a) and L8, even with full cardinality coverage]
\label{C:K-counterexample}
Axioms a)--d) of cote~68, scan~017, do not imply Condition~K of scan~021. This remains false when $\rho$ is surjective, its kernel is exactly the specified copy of $G[2]$, every chain is nonempty, every chain contains $G$, every singleton occurs, and every cardinality from $1$ to $|G|$ occurs.
\end{theorem}
\begin{proof}
Take $G=\mathbb Z/4\mathbb Z$, $\widetilde G=\mathbb Z/8\mathbb Z$,
$\rho(k)=k\pmod4$, and $i(2)=4$. For $t\in\mathbb Z/4\mathbb Z$ define
\begin{align*}
\mathcal A(2t)&=\bigl\{\{t\},\{t-1,t,t+1\},G\bigr\},\\
\mathcal A(2t+1)&=\bigl\{\{t,t+1\},G\bigr\}.
\end{align*}
Each is a chain and its sets are invariant under $x\mapsto k-x$ for the corresponding label $k$. Adding $4$ to the label translates the chain by $2$, verifying the kernel covariance. Translate two chains so their centers are respectively $0$ and either $0$ or $1/2$. In the mixed case their sets interlace as
$\{0\}\subset\{0,1\}\subset\{-1,0,1\}\subset G$;
thus the final nesting axiom holds.

It remains to verify convolution closure, which can be reduced by translation to three elementary identities. Write $D=1_{\{0\}}$, $E=1_{\{-1,0,1\}}$, $F=1_{\{0,1\}}$, and $\mathbf1=1_G$. Then
\[
E*E=D+2\mathbf1,\qquad
E*F=F+\mathbf1,\qquad
F*F=L_1(D+E).
\]
Convolution with $D$ is the identity, and $1_A*\mathbf1=|A|\mathbf1$. These formulas cover every generator pair and have nonnegative coefficients in the required output chain; the output labels are respectively $0,1,2$. This proves every axiom, not merely the peak computation.

Now take all three input functions equal to $f=F$, and all their rearrangements equal to $f'=1_{\{0,2\}}$. In the order $0,1,2,3$,
\[
f*f*f=(1,3,3,1),\qquad f'*f'*f'=(4,0,4,0).
\]
Thus the asserted peak inequality would read $4\leq3$. Even two-factor majorization fails: $f*f=(1,2,1,0)$ has its two largest entries summing to $3$, whereas $f'*f'=(2,0,2,0)$ has sum $4$.

The attached exact certificate checks all eight chains, twenty reflection identities, sixteen kernel-covariance identities, sixty-four chain alignments, and all four hundred labelled generator convolutions. It also verifies all four cardinalities and both displayed convolution vectors. No floating-point approximation is used.
\end{proof}

\begin{theorem}[Proved: an exact additional condition that repairs L8]
\label{C:K-repair}
Assume the system axioms and that $S$ contains an indicator of every possible finite cardinality. Then the following are equivalent:
\begin{enumerate}
\item every pair of generating indicators, from any two labelled chains, is tass\'e;
\item every pair in $S^2$ is tass\'e;
\item for every $n\geq2$, $f_i\in S$ and $f_i'\preceq f_i$ imply
$f_1'*\cdots*f_n'\preceq f_1*\cdots*f_n$;
\item Condition~K: for the same data,
$\|f_1'*\cdots*f_n'\|_\infty\leq\|f_1*\cdots*f_n\|_\infty$.
\end{enumerate}
The cardinality assumption is needed only for (4)$\Rightarrow$(2); (1)$\Leftrightarrow$(2)$\Leftrightarrow$(3)$\Rightarrow$(4) holds without it. Thus a)--d) provide compatibility of the proposed extremizers, while extremality of generator pairs is an additional mathematical condition.
\end{theorem}
\begin{proof}
Only (1)$\Rightarrow$(2) and (4)$\Rightarrow$(2) require proof. Write
$f=\sum_i\alpha_i1_{A_i}\in C(\tilde a)$ and
$g=\sum_j\beta_j1_{B_j}\in C(\tilde b)$ with nested sets and positive coefficients. A rearrangement $f'\sim f$ transports the whole level chain to a chain $A_i'$ of the same cardinalities; similarly for $g'$. By (1),
$1_{A_i'}*1_{B_j'}\preceq1_{A_i}*1_{B_j}$.
All right-hand functions lie in the same cone $C(\tilde a\tilde b)$, so one decreasing order works for all of them. The subadditivity of $p_k$ and its additivity on such simultaneously ordered functions give
\begin{align*}
p_k(f'*g')
&\leq\sum_{i,j}\alpha_i\beta_jp_k(1_{A_i'}*1_{B_j'})\\
&\leq\sum_{i,j}\alpha_i\beta_jp_k(1_{A_i}*1_{B_j})
 =p_k(f*g).
\end{align*}
Lemma~\ref{C:majorization} extends this from rearrangements to weakly majorized inputs by separate convexity and positivity. This proves (2).

For (4)$\Rightarrow$(2), fix $k$ and choose $1_H\in S$ with $|H|=k$. Choose a $k$-point set $H'$ so that
$(f'*g'*1_{H'})(e)=p_k(f'*g')$; one takes inverses of points realizing the $k$ largest values, with arbitrary zero-value points added if needed. Applying (4) to three factors gives
\[
p_k(f'*g')\leq\|f*g*1_H\|_\infty=p_k(f*g).
\]
The equality is Theorem~\ref{C:pretasse-cones}, since $f*g\in S$. This proves (2). Finally (2) iterates using convolution closure to give (3), and taking $p_1$ gives (4). The reverse implication (2)$\Rightarrow$(1) is immediate.
\end{proof}

\subsection{Which groups carry the printed system, and what the kernel means}
The displayed data on scan~017 include a group isomorphism $i:G[2]\simeq\ker\rho$, and one of the lettered axioms uses $i$ in its translation covariance. We therefore retain that covariance throughout. On the standard reading $G[2]=\{u\in G:u^2=e\}$, the full displayed system has the following classification. We also give the alternative answer if the notation denotes a chosen subgroup of exponent two. The lettered axioms do not impose cardinality coverage, nonempty fibers, or faithfulness of the indexing group.

\begin{theorem}[Proved under the explicit reading: literal classification, C.2(b)]
\label{C:group-classification}
Suppose $G[2]$ means the full set of elements whose square is $e$, and the displayed $i:G[2]\simeq\ker\rho$ is a group isomorphism. With finite subsets and empty families allowed, a nonzero system satisfying all the displayed data and all lettered axioms exists if and only if the elements of $G[2]$ commute pairwise. Surjectivity of $\rho$ does not change this answer. For finite groups satisfying this condition, all fibers can also be required nonempty.

If instead $G[2]$ is interpreted as a specified subgroup $H\leq G$ of exponent two, every ambient group $G$ admits such a system with $\ker\rho\simeq H$. This includes the choice $H=\{e\}$ for every group. In either interpretation, the requirement that $i$ map \emph{onto} $\ker\rho$ is not forced by the chain, reflection, convolution, alignment, and kernel-covariance axioms, even with full finite cardinality coverage and $i$ still injective into the kernel.
\end{theorem}
\begin{proof}
The full set $G[2]$ is a subgroup if and only if its elements commute pairwise. Indeed, if it is a subgroup, then for $u,v\in G[2]$ one has $uv=(uv)^{-1}=vu$; conversely commuting involutions have product an involution. The displayed group-isomorphism condition presupposes this subgroup property.

For sufficiency, put $H=G[2]$, or use the specified subgroup $H$ under the alternative interpretation. Such an $H$ is elementary abelian. Set $\widetilde G=H\times G$, $\rho(h,g)=g$, and $i(u)=(u,e)$. Define
\[
\mathcal A(h,e)=\{\{h\}\},\qquad
\mathcal A(h,g)=\varnothing\quad(g\ne e).
\]
This is a nonzero system, $\rho$ is surjective, and $i$ identifies $H$ with its exact kernel. For each nonempty family the reflection equation is $\{h\}=e\{h\}^{-1}$, because $h^{-1}=h$. The only nonzero generator convolutions are $1_{\{h\}}*1_{\{k\}}=1_{\{hk\}}$, in the correctly labelled output family. Translation by $u\in H$ changes $\{h\}$ to $\{uh\}$, exactly $\mathcal A(i(u)(h,e))$; it preserves empty families as required at other labels. Any two singleton families align by a left translation, and joint nesting with an empty family is automatic. Thus every axiom, including covariance, holds. For finite $G$, one may instead set every family to $\{G\}$: reflection, covariance, and alignment are immediate, and $1_G*1_G=|G|1_G$. This gives nonempty fibers.

Finally, given any system, replace its indexing group by $\widetilde G\times Q$ for any group $Q$, compose $\rho$ with projection, set $i'(u)=(i(u),e_Q)$, and repeat the old family at every new label. All lettered axioms, including covariance through $i'$, and any coverage property survive. The kernel becomes $(\ker\rho)\times Q$, strictly larger than $i'(H)$ when $Q$ is nontrivial. Apply this to the full-coverage $C_4$ example with $Q=C_3$ to obtain a nondegenerate finite example whose kernel strictly exceeds $G[2]$. Thus exactness at the kernel is additional data, not a consequence of the other axioms.
\end{proof}

\begin{proposition}[Proved: the integer interval system]
\label{C:integer-system}
For $G=\mathbb Z$, a nonempty system representing every finite cardinality is obtained with $\widetilde G=\mathbb Z$, $\rho=1$, and $\mathcal A(k)$ consisting of all finite integer intervals whose midpoint is $k/2$. Its kernel is $0=G[2]$. The same chain, reflection, and convolution construction descends to cyclic groups of even order $N$ with indexing group $\mathbb Z/(2N)$, $\rho(k)=k\pmod N$, and $i(N/2)=N$, after including the whole group in every chain.
\end{proposition}
\begin{proof}
On $\mathbb Z$, the admissible interval lengths have parity opposite to $k$. The intervals about a fixed midpoint are nested and symmetric under $x\mapsto k-x$. Convolution of intervals is the overlap-count function of two intervals: its values increase to a plateau and decrease symmetrically, and its positive level sets are intervals about the sum of the midpoints. Therefore it is a nonnegative sum of indicators in the required output chain. Two chains can be translated to centers $0$ and $0$ or $1/2$, where the intervals interlace. This proves every assertion on $\mathbb Z$.

For the cyclic statement, an interval of length $a\leq N$ and midpoint $k/2$ is represented by consecutive residues with endpoints $(k-a+1)/2$ and $(k+a-1)/2$ when these are integers; the full group is added separately. Translating a midpoint by $N/2$ translates the set by the unique nonzero element of $G[2]$, which proves kernel covariance. To check convolution closure, first suppose the lengths $a,b$ satisfy $a+b\leq N+1$. On $\mathbb Z$, the convolution of $\{0,\ldots,a-1\}$ and $\{0,\ldots,b-1\}$ has values
$\max(0,\min(t+1,a,b,a+b-1-t))$ at $t$. It is symmetric about $(a+b-2)/2$ and has nested interval level sets. Its support has length at most $N$, so reduction modulo $N$ preserves that layer decomposition. If $a+b>N+1$ and neither set is the whole group, use the exact identity
\[
1_A*1_B=(a+b-N)1_G+1_{G\setminus A}*1_{G\setminus B}.
\]
The complements are intervals of lengths $N-a,N-b$ whose sum is less than $N-1$, so the preceding case applies. Each complement has its midpoint shifted by $N/2$, and the two shifts sum to $N$, giving the same output midpoint modulo $N$. The constant term is nonnegative. If either factor is $G$, the convolution is constant. Thus every positive level set is a centered cyclic interval or the whole group. This proves the required nonnegative layer decomposition. Joint nesting follows as in the integer case. The exact $N=4$ certificate supplies an independent finite verification of the case used in Theorem~\ref{C:K-counterexample}.
\end{proof}

\begin{proposition}[Proved: an exact-kernel obstruction already on the odd cycle $C_3$]
\label{C:C3-coverage-obstruction}
Let $G=\mathbb Z/3\mathbb Z$ and retain the printed exact-kernel condition $\ker\rho=G[2]=0$. Any system satisfying the chain and convolution axioms has no two-element generating set in any family. Consequently it cannot have full cardinality coverage. Surjectivity of $\rho$ and the remaining printed axioms do not remove this obstruction.
\end{proposition}
\begin{proof}
Because $\rho$ is injective into $C_3$, every label $h\in\widetilde G$ satisfies $h^3=e$. Suppose $A=G\setminus\{m\}$ belongs to the family at $h$. Write $\mathbf1=1_G$ and $\delta_j=1_{\{j\}}$. Exact convolution gives
\[
 1_A*1_A=\mathbf1+\delta_{2m},\qquad
 1_A*1_A*1_A=2\mathbf1+1_{G\setminus\{0\}}.
\]
In a cone generated by a chain of finite sets, every nonempty positive level set of every cone element is a member of that chain: write the function as a nonnegative combination of distinct nested indicators, discard zero coefficients, and its successive positive levels are exactly the sets that remain. The first identity and convolution closure therefore force $\{2m\}$ into the family at $h^2$. The second identity forces $G\setminus\{0\}$ into the family at $h^3=e$. Cubing the singleton generator gives $\delta_{2m}^{*3}=\delta_0$ in the family at $h^6=e$, so that family also contains $\{0\}$. These two nonempty sets are incomparable, contradicting the chain axiom.

The certificate \texttt{code/C\_C3\_impossibility.py} independently exhausts all $12^3=1728$ systems of reflection-compatible chains at the three possible labels, allowing empty families and the empty set. None of the convolution-closed systems contains a two-element set. Since an injective indexing group is a subgroup of $C_3$, the case with only the identity label embeds in this check by leaving the other families empty; the proof itself does not rely on this enumeration.
\end{proof}

\begin{remark}[Proved: scope of the classification]
\label{C:classification-scope}
Theorem~\ref{C:group-classification} classifies the \emph{printed} existence problem under each explicitly stated interpretation of its two-torsion notation. Requiring all cardinalities in every infinite system, all fibers nonempty for an infinite group, or a minimal faithful indexing group specifies different problems. We have not silently claimed that these strengthened classifications follow from the printed axioms. The interval construction above is deliberately stated for even cyclic groups: on an odd cycle, the doubled midpoint index has a two-element kernel although $G[2]=0$, so it does not satisfy the printed exact-kernel datum. Proposition~\ref{C:C3-coverage-obstruction} shows that a full-cardinality exact-kernel substitute is impossible already for $C_3$. This does not contradict the unrestricted finite existence construction using only $\{G\}$. In particular the kernel notation for a genuinely nonabelian group must be typed: if its involutions do not commute, $G[2]$ is not a group and the displayed group isomorphism does not exist. This is separate from the rearrangement definitions, which make sense on every group.
\end{remark}

\subsection{Symmetry: a complete noncommutative counterexample}
\begin{theorem}[Refuted in general, with exact replacement: C.2(c)]
\label{C:symmetry}
Neither pr\'etass\'e nor tass\'e is a symmetric relation on arbitrary groups. Both are symmetric on abelian groups. On every group their universally valid reversal is
\[
(f,g)\text{ has the property}
\quad\Longleftrightarrow\quad
(\check g,\check f)\text{ has the property}.
\]
For the pr\'etass\'e property, symmetry also holds whenever each function is invariant under inversion followed by some left translation.
\end{theorem}
\begin{proof}
For the counterexample take $G=S_3$, $H=\{1,(12)\}$ and $s=(123)$. Set $f=1_H$, $g=1_{Hs}$. Since $1_H*1_H=2\,1_H$,
$f*g=2\,1_{Hs}$ and its peak is $2=R(f,g)$. On the other hand
$g*f=1_H*\delta_s*1_H$ is the indicator of $HsH$, a four-element set: the four products $hsk$ are distinct because $H\cap sHs^{-1}=\{1\}$. Hence its peak is $1$, so $(g,f)$ is not pr\'etass\'e.

In fact $(f,g)$ is tass\'e. Any $f'\preceq f$, $g'\preceq g$ has entries at most one and total mass at most two. Thus each entry of $f'*g'$ is at most two and its total mass is at most four. It follows that
$p_k(f'*g')\leq\min(2k,4)=p_k(f*g)$.
The reversed pair cannot be tass\'e because it is not pr\'etass\'e. The attached program independently checks all $15^2=225$ pairs of two-point subsets in $S_3$, as well as the exact two convolution vectors.

For the general reversal, inversion permutes values and
$\check{(f*g)}=\check g*\check f$; it therefore preserves both the peak and all decreasing partial sums, and sends every rearrangement or weak-majorization comparison to one of the same type. On an abelian group $f*g=g*f$, which proves symmetry directly.

Finally suppose $\check f=L_a f$ and $\check g=L_b g$ for some $a,b\in G$. By Lemma~\ref{C:majorization}, $(f,g)$ is pr\'etass\'e precisely when the level sets of $f$ and of some left translate of $g$ are jointly nested. Translating this nesting by its inverse is equivalent to nesting $g$ with a left translate of $f$, which is precisely the reverse pr\'etass\'e condition. We make no analogous assertion for tass\'e pairs without further hypotheses.
\end{proof}

\paragraph{Opened HLP source and its precise scope.}
The copy actually opened is the original 1934 edition of
\emph{Inequalities}, not the unverified 1952 second edition.
Chapter~X, Theorems~368--370 (\S\S10.2--10.4, pp.~261--265),
gives the finite scalar-product rearrangement and two-factor
coefficient comparisons corresponding to the rearrangement part of
Lemma~\ref{C:majorization}.  For its additional convolution comparison,
write \(a^+\) for the values of a finitely supported nonnegative
sequence on \(\mathbb Z\), placed in decreasing order at
\(0,1,-1,2,-2,\ldots\), and put \(a^-(j)=a^+(-j)\).
The source calls a nonzero multiset symmetric when its largest
positive value has odd multiplicity and its other positive values
have even multiplicities; zeros may be added to obtain the book's
odd-length indexing interval.  In that case \(a^+=a^-=a^\star\)
is even and decreases with \(|j|\).
With \(a,b\) arbitrary and each of \(c,d,\ldots\) symmetric in
this sense, Theorem~373 and its extension, Theorem~374, state
\[
 (a*b*c*d*\cdots)(0)
 \ \leq\
 (a^+*b^-*c^\star*d^\star*\cdots)(0).
\]
The complete finite proof is in \S\S10.5--10.9, pp.~265--274:
it reduces to indicator layers, proves the three-factor comparison,
and uses closure of symmetric decreasing convolution
(Theorem~375, p.~273) for the induction.
These exact source hypotheses are retained here
\cite[\S\S10.1--10.9, pp.~260--274]{C_HLP1934}.

This source comparison concerns rearrangements of fixed finite
multisets.  The separate positivity and convex reduction in
Lemma~\ref{C:majorization} is the justification used in this
manuscript for passing from suitable rearrangement inequalities to
weak majorization; one cannot instead assume that an arbitrary
weakly majorized input retains the displayed symmetry.
Nor can the integer coefficient relation in HLP be transferred
to an arbitrary group: Theorem~\ref{C:K-counterexample} gives an
explicit obstruction to that proposed Condition~K implication.
HLP \S10.10 itself retains only a multiplicative constant depending
on the number of factors in its unrestricted Theorem~377
(pp.~274--276), rather than the same prescribed-order
constant-one comparison.

Finally, these inequalities do not construct the labelled chains,
the map \(\rho\), its specified kernel, or the covariance identities
of the manuscript's system.  Proposition~\ref{C:integer-system}
checks those separately, including intervals of even cardinality,
which have half-integer centres.  An even-cardinality indicator
does not satisfy the symmetric-multiset hypothesis above, and
Theorem~374 allows only two arbitrary multisets.  Thus these quoted
HLP results alone do not establish Condition~K for the full integer
system with arbitrarily many such inputs.  This is a limitation of
the cited comparison, not an assertion that Condition~K fails for
that integer system.  The original Gabriel article has now also
been opened: Theorems~1--2 (pp.~32--39) give the finite three-factor
and additional-symmetric-factor comparisons above, while
Theorems~3--4 (pp.~39--42) concern paired autocorrelations and
products of squared moduli of finite trigonometric polynomials
\cite{C_Gabriel1932}. Their paired structure does not supply a
comparison for unrestricted independent factors, and none of these
results supplies the manuscript's additional system axioms.


\paragraph{Literature boundary.}
We searched for Grothendieck's cote~68 rearrangement terminology, arbitrary module-functor exponential maps, discrete group rearrangement systems, and the Gabriel and Lev references. We opened the primary arXiv paper of Madiman, Wang, and Woo (2020 version; published 2021)~\cite{C_MWW2021}, including its source file, and checked its Definitions~2.1--2.3 and Theorems~6.2 and~6.8. Its cyclic rearrangement results concern \emph{prime} cyclic groups. They cannot justify the all-cyclic-group extrapolation in the attached study; Theorem~\ref{C:K-counterexample} rules that extrapolation out. No priority is claimed for our elementary arguments or finite counterexamples. Bibliographic entries record whether the actual source was opened; an unsuccessful download is not recorded as an opened theorem.

\subsection{Graded determinants: the requested source and its qualifications}
% Notes-only proposal. Root owns integration, numbering and register changes.
% The three corrections are auxiliary source audits, not movement of either
% outstanding subject-21 original-source gate.
\paragraph{Requested determinant background: conventions and scope.}
The original of cote~68, scan~056, concerns finite free modules,
free rank-one graded modules with Koszul commutativity, and invariance
under changing a free complement.  It is a finite-free prototype of the
determinant formalism.  The bounded-complex and perfect-complex data in
Knudsen--Mumford's Definitions~1 and~4 have additional requirements
\cite[pp.~20--24, 40--42]{ROOT_KnudsenMumford1976}.
The following proofs also record three qualifications found by reading
all of the original article, pp.~19--55.

\begin{lemma}[Proved: graded signs, exterior blocks and the printed sum-exponent defect]\label{C:graded-determinant-sign}
Let $S$ be any scheme.  Let $\operatorname{Pic}^{\mathbb Z}(S)$ be the
groupoid whose objects are pairs $(L,a)$, with $L$ an invertible
$\mathcal O_S$-module and $a:S\to\mathbb Z$ locally constant, and whose
arrows are degree-preserving $\mathcal O_S$-linear isomorphisms.  Set
\[
 (L,a)\otimes(M,b)=(L\otimes M,a+b),\qquad
 \symbf{1}=(\mathcal O_S,0).
\]
With the ordinary tensor associator and unit maps, the formula
\[
 c_{(L,a),(M,b)}(l\otimes m)=(-1)^{ab}m\otimes l
\]
defines a symmetric monoidal Picard groupoid.  The functor
\[
 E\longmapsto\operatorname{Det}(E)
       =(\textstyle\bigwedge^{\operatorname{rk}E}E,\operatorname{rk}E)
\]
on finite locally free modules and their isomorphisms, equipped with
the exterior-product maps for direct sums, respects this symmetry.
It respects all finite block permutations and, on finite free modules,
the change-of-complement condition of scan~056.

The replacement $(-1)^{a+b}$ printed in the commutativity formula of
Knudsen--Mumford, p.~20, fails these unit and exterior-product
requirements already over $\operatorname{Spec}\mathbb Q$.
\end{lemma}
\begin{proof}
The function $(-1)^{ab}$ is a locally constant unit of
$\mathcal O_S$, so the displayed formula is a globally defined line
isomorphism.  It is natural because line morphisms are linear.
The composite of the two swaps has scalar
$(-1)^{ab+ba}=1$.  Under the ordinary tensor associator the two
hexagon identities reduce, respectively, to
\[
 a(b+c)=ab+ac,\qquad (a+b)c=ac+bc.
\]
The unit compatibility follows from $a\cdot0=0$.
The associator has its ordinary tensor pentagon, and
$(L^\vee,-a)$ is a tensor inverse of $(L,a)$, by evaluation.
This verifies the stated Picard structure for arbitrary integral
degrees, including negative ones.

All assertions about finite locally free modules can be checked where
their ranks are constant and bases exist.  For bases
$e_1,\ldots,e_r$ of $E$ and $f_1,\ldots,f_s$ of $F$, let
$i_{E,F}$ send the tensor of the two volume forms to their concatenated
wedge in $\bigwedge^{r+s}(E\oplus F)$.  Let
$\tau:E\oplus F\to F\oplus E$ send $(e,f)$ to $(f,e)$.
Moving the $r$ vectors of $E$ past the $s$ vectors of $F$ makes exactly
$rs$ elementary exterior swaps.  Consequently
\[
 \operatorname{Det}(\tau)\,i_{E,F}
       =i_{F,E}\,c_{\operatorname{Det}(E),\operatorname{Det}(F)}.
\]
Associativity of the maps $i$ is associativity of concatenation;
their unit normalization uses the empty wedge.  For a permutation
which sends block $i$ to position $\sigma(i)$, every inverted pair
$i<j$ with $\sigma(i)>\sigma(j)$ contributes $r_i r_j$ crossings.
Both the exterior and graded-tensor routes therefore have the sign
\[
 (-1)^{\sum_{i<j,\ \sigma(i)>\sigma(j)}r_i r_j}.
\]
This proves the assertion for every finite block permutation.

For the manuscript's complement condition, write $V=W\oplus U$
and $V=W\oplus U'$ with all indicated modules finite free.
Let $\lambda:U\to U'$ be the isomorphism inducing the identity
on $V/W$.  For a basis $w_1,\ldots,w_r$ of $W$ and a basis
$u_1,\ldots,u_s$ of $U$, put $u'_j=\lambda(u_j)$.
Then $u'_j-u_j\in W$.  In
\[
 w_1\wedge\cdots\wedge w_r\wedge u'_1\wedge\cdots\wedge u'_s
\]
every expansion term other than the one with all the $u_j$ vanishes,
since it contains $r+1$ vectors from $W$.  The two determinant
identifications are equal.  Equivalently, an automorphism which is
the identity on $W$ and $V/W$ has block form
$\left(\begin{smallmatrix}I&T\\0&I\end{smallmatrix}\right)$,
and the same wedge calculation gives determinant one.  The rank-one
normalization is $\bigwedge^1L=L$ in degree one.

Finally, over $\mathbb Q$, degrees $(1,0)$ give scalar $-1$ under
the printed sum rule, although swapping with the specified unit has
scalar $+1$.  For two rank-one blocks, degrees $(1,1)$ give $+1$
under that rule, whereas $f\wedge e=-e\wedge f$.  Thus both are
actual counterexamples in characteristic zero.  The sum rule is
nevertheless involutive, so the double-swap identity alone would not
detect its defect.
\end{proof}

Grothendieck's letter to Knudsen dated 19 May 1973, as reproduced in
the opened Appendix~B of Knudsen's later paper, explicitly uses the
product of the two degrees and requires associativity and unit
compatibility \cite[Appendix~B, author-version pp.~40--42, especially
p.~41]{A567_GrothendieckKnudsen1973Letter}.
This supplies original-context corroboration; it is not a claim to
have located a formal erratum to the 1976 printing.
The exact finite certificate \texttt{C\_graded\_determinant\_sign\_certificate.py}
checks four parity pairs, eight parity triples and the 25 block-rank
pairs with ranks between zero and four.  Its 90 assertions passed.
The general assertions are established by the proof above, not by
extrapolation from those finite cases.

\begin{proposition}[Refuted: unrestricted ordinary base change in the printed pushforward paragraph]\label{C:determinant-ordinary-basechange}
The assumptions ``$p:X\to Y$ proper of finite Tor dimension,
$Y$ noetherian, $F$ perfect on $X$, and $g:Y'\to Y$ a morphism
of noetherian schemes'' do not imply the ordinary-fibre-product
comparison displayed on p.~46 of
\cite{ROOT_KnudsenMumford1976}.  With
$X'=X\times_Y Y'$ and projections $g':X'\to X$, $p':X'\to Y'$,
the two graded lines
\[
 \operatorname{Det}(Lg^*Rp_*F)
 \quad\hbox{and}\quad
 \operatorname{Det}(Rp'_*Lg'^*F)
\]
need not even have the same degree.
\end{proposition}
\begin{proof}
Put $R=\mathbb Q[t]$, $k=R/(t)$,
$Y=\operatorname{Spec}R$ and
$X=Y'=\operatorname{Spec}k$.  Let both $p$ and $g$ be the closed
immersion induced by $R\to k$, and set $F=\mathcal O_X$ in degree
zero.  The morphism $p$ is proper, and all schemes are noetherian.
The exact free resolution
\[
 0\longrightarrow R\xrightarrow{\ t\ }R\longrightarrow k
   \longrightarrow0
\]
shows that $p$ has finite Tor dimension one.  The complex $F$ is
perfect on $X$.  The ordinary tensor product $k\otimes_R k=k$
identifies $X'$ with $\operatorname{Spec}k$; both $g'$ and $p'$
are identity maps.

Since a closed immersion is affine, $Rp_*F$ is the $R$-module $k$
in degree zero.  Pulling its displayed free resolution back gives
\[
 Lg^*Rp_*F=k\otimes_R^{\mathbf L}k
       = [\,k\xrightarrow{\ 0\ }k\,],
       \qquad\text{in degrees }-1,0.
\]
Its graded determinant has degree $-1+1=0$.  In contrast,
$Rp'_*Lg'^*F=k$ in degree zero has determinant degree one.
Degree-preserving line isomorphisms cannot identify these objects.
The derived comparison is already not a quasi-isomorphism, since
its source has $H^{-1}=k$ and its target has $H^{-1}=0$.
The obstruction is the nonzero
$\operatorname{Tor}_1^R(k,k)=k$; the square is not Tor independent.
The first isomorphism printed on p.~46, between the pullback of a
determinant line and the determinant of the derived pullback, is not
challenged by this example.  It is the second, pushforward base-change
comparison that fails.
\end{proof}

\begin{proposition}[Proved: the sharp affine criterion and Tor-independent repair]\label{C:determinant-basechange-repair}
Let $R\to S$ and $R\to A$ be homomorphisms of commutative rings,
and put $B=A\otimes_R S$.  For a perfect $S$-complex $F$, consider
the classical comparison in $D(A)$
\[
 \beta_F:A\otimes_R^{\mathbf L}F
       \longrightarrow B\otimes_S^{\mathbf L}F,
\]
where the target is restricted along $A\to B$.
Then $\beta_F$ is a quasi-isomorphism for every perfect $S$-complex
$F$ if and only if $\operatorname{Tor}_i^R(A,S)=0$ for all $i>0$.

More generally, let $p:X\to Y$ be a quasi-compact, quasi-separated
morphism of schemes, let $g:Y'\to Y$ be any morphism, and form
$X'=X\times_Y Y'$ with projections $g'$ and $p'$.  Suppose that
$X$ and $Y'$ are Tor independent over $Y$: on affine charts over
an affine open of $Y$, all higher Tor modules of their two coordinate
algebras vanish.  For every
$E\in D_{\mathrm{QCoh}}(\mathcal O_X)$ the canonical map
\[
 Lg^*Rp_*E\longrightarrow Rp'_*Lg'^*E
\]
is an isomorphism.  If $Rp_*E$ is perfect, this gives an isomorphism
of their graded determinants as well.  In particular, flatness of
either $p$ or $g$ supplies the Tor-independence hypothesis.
When $p$ is proper of finite Tor dimension, $Y$ is noetherian and
$E$ is perfect, $Rp_*E$ is perfect, so this is a corrected version
of the printed determinant comparison.
\end{proposition}
\begin{proof}
For necessity in the affine assertion take $F=S[0]$.
The target is $B[0]$, whereas the source has
$H^{-i}=\operatorname{Tor}_i^R(A,S)$ for $i\geq0$.
For sufficiency choose a free $R$-module resolution $P^\bullet\to A$
in nonpositive degrees and represent $F$ by a bounded complex
$G^\bullet$ of finite projective $S$-modules.  After forgetting the
$A$-action, the canonical comparison is represented by
\[
 (P^\bullet\otimes_R S)\otimes_S G^\bullet
       \longrightarrow B\otimes_S G^\bullet.
\]
The assumed Tor vanishing makes $P^\bullet\otimes_R S\to B$
a quasi-isomorphism of $S$-complexes.  Tensoring with a single
projective $S$-module preserves an acyclic complex.  The finite
filtration of $G^\bullet$ by its degrees therefore shows, successively,
that tensoring with $G^\bullet$ preserves this quasi-isomorphism.
Forgetting from $D(A)$ to $D(R)$ detects quasi-isomorphisms, since
it does not change the underlying cohomology groups.  This proves
sufficiency without incorrectly assigning an $A$-module structure
to an arbitrary chosen $R$-free resolution.

For the scheme assertion use the canonical base-change comparison
and cover $Y'$ by affine opens lying over affine opens of $Y$.
Restriction to these opens commutes with the sheaf-theoretic
derived images in question.  It suffices to take $g$ between affine schemes.
Then $g$ and $g'$ are affine, and $Rg_*$ is conservative on complexes
with quasi-coherent cohomology.  Indeed, affine pushforward is exact
and faithful on quasi-coherent modules, so it detects each cohomology
module.  Write $K=g_*\mathcal O_{Y'}$.  Pushing the comparison forward
along $g$ identifies it with
\[
 Rp_*E\otimes_{\mathcal O_Y}^{\mathbf L}K
       \longrightarrow
 Rp_*\bigl(E\otimes_{\mathcal O_X}^{\mathbf L}
                          g'_*\mathcal O_{X'}\bigr).
\]
The ordinary affine fibre-product calculation gives
$g'_*\mathcal O_{X'}=p^*K$.  Tor independence gives
$Lp^*K\simeq p^*K$, as follows on an affine chart from the
vanishing of the higher homology of the derived tensor product.
The displayed map is thus the derived projection-formula
isomorphism for the quasi-compact, quasi-separated map $p$.
This is the derivation in the opened Stacks proofs
\cite[Tags~08IB and 08EU]{BStacks}.
The latter projection formula can be checked after taking $Y$
affine: its valid objects $K$ are closed under shifts, triangles
and arbitrary direct sums, since derived tensor and $Rp_*$ on
quasi-coherent complexes preserve those sums.  Free resolutions
and their truncation homotopy colimits generate $D(R)$ from $R$,
where the formula is the identity.  This proves it for the required
$K$ and completes the comparison.

Derived pullback preserves a bounded locally free local model, so
if $Rp_*E$ is perfect then the left-hand object is perfect; the
isomorphism makes the right-hand object perfect too.  Their
determinants may therefore be taken.  If one side of the square
is flat, its local algebra is flat over the base, so its higher
Tor modules with the other algebra vanish.

For the final proper case, work locally on the noetherian target.
A perfect complex on the proper source has bounded coherent
cohomology and finite Tor amplitude relative to the target:
use its bounded locally free models and the finite Tor dimension
of $p$, with a finite cover over the chosen affine base.
Proper coherent pushforward gives bounded coherent cohomology
for $Rp_*E$, hence pseudo-coherence.  For every quasi-coherent
module $M$ on the base, the projection formula identifies
$Rp_*E\otimes^{\mathbf L}M$ with
$Rp_*(E\otimes_{p^{-1}\mathcal O_Y}^{\mathbf L}p^{-1}M)$.
Relative Tor amplitude bounds the latter uniformly from below;
the bounded-above complex $Rp_*E$ supplies the upper bound.
Thus $Rp_*E$ has finite Tor amplitude.  Pseudo-coherence and finite
Tor amplitude imply perfection.  These are precisely the steps
in the opened perfect-pushforward proof
\cite[Tags~0B6G and 08EV]{BStacks}.
\end{proof}

\begin{lemma}[Proved: the exact scalar trace-splitting condition]\label{C:determinant-trace-splitting}
Let $R$ be a commutative ring, let $m\geq1$, and set
$\mathfrak s_m(R)=\ker(\operatorname{tr}:M_m(R)\to R)$.
The natural map
\[
 \phi:R\oplus\mathfrak s_m(R)\longrightarrow M_m(R),
       \qquad (a,T)\longmapsto aI_m+T,
\]
fits into the exact sequence
\[
 0\longrightarrow R[m]\longrightarrow
 R\oplus\mathfrak s_m(R)\xrightarrow{\phi} M_m(R)
 \xrightarrow{\operatorname{tr}\bmod m}R/mR\longrightarrow0,
\]
where $R[m]=\{a\in R:ma=0\}$ and the first map sends
$a$ to $(a,-aI_m)$.
It is an isomorphism if and only if $m$ is a unit in $R$.
In particular, the scalar/trace-zero direct-sum description on p.~52
of \cite{ROOT_KnudsenMumford1976} is false over the unrestricted
noetherian bases stated there.  For $R=\mathbb F_2$, $m=2$ there
is not even a lift of $1$ under trace which is invariant under
every change of basis.
\end{lemma}
\begin{proof}
The equation $aI_m+T=0$ forces $T=-aI_m$, and this has trace zero
exactly when $ma=0$.  This proves the kernel assertion and the
injectivity of the first map.  The final trace map is onto because
$\operatorname{tr}(E_{11})=1$.  A matrix is in its kernel exactly
when its trace is $ma$ for some $a\in R$; subtracting $aI_m$ then
makes it trace zero.  This proves exactness at $M_m(R)$.

If $m$ is a unit, the inverse of $\phi$ is
\[
 A\longmapsto
 \bigl(m^{-1}\operatorname{tr}(A),\,
       A-m^{-1}\operatorname{tr}(A)I_m\bigr).
\]
Conversely, if $\phi$ is surjective, express $E_{11}=aI_m+T$
with $\operatorname{tr}(T)=0$.  Taking traces gives $1=ma$,
so $m$ is a unit.  This proves the equivalence.

For the explicit noetherian example $R=\mathbb F_2$, $m=2$,
the nonzero identity matrix has trace zero.  Thus the scalar line
intersects the trace-zero subspace nontrivially; equivalently
$\phi(1,I_2)=0$ with $(1,I_2)\ne0$.

For the stronger invariance assertion put
\[
 J=\begin{pmatrix}0&1\\1&0\end{pmatrix},\qquad
 U=\begin{pmatrix}1&1\\0&1\end{pmatrix}.
\]
Both matrices are invertible over $\mathbb F_2$.
Writing $M=\left(\begin{smallmatrix}a&b\\c&d\end{smallmatrix}\right)$,
the equation $MJ=JM$ gives $c=b$ and $d=a$.
Substitution in $MU=UM$ then gives $b=0$.  Thus a matrix fixed by
conjugation by both $J$ and $U$ is scalar, and its trace is
$2a=0$.  Any change-of-basis-invariant lift of $1$ would commute
with both matrices and have trace one, a contradiction.
\end{proof}

\begin{lemma}[Proved: incidence and its Koszul resolution without trace splitting]\label{C:determinant-incidence}
Let $E$ be a finite locally free module of positive constant rank
$m$ on any scheme $S$.  Use the convention that $\mathbf P(E)$
parametrizes rank-one locally free quotients of $E$.
On $W=\mathbf P(E)\times_S\mathbf P(E^\vee)$ there is a canonical
section
\[
 \delta\in\Gamma(W,
  \operatorname{pr}_1^*\mathcal O_{\mathbf P(E)}(1)
  \otimes\operatorname{pr}_2^*\mathcal O_{\mathbf P(E^\vee)}(1))
\]
whose zero scheme is the incidence locus.  It commutes with arbitrary
base change, and
\[
 0\longrightarrow\mathcal O_W(-1,-1)
       \xrightarrow{\ \delta\ }\mathcal O_W
       \longrightarrow\mathcal O_{V(\delta)}\longrightarrow0
\]
is exact.  On a product with finitely many independent dual factors,
the tensor product of these two-term complexes resolves their
scheme-theoretic incidence intersection.  These assertions need
no invertibility condition on $m$.
\end{lemma}
\begin{proof}
For any finite locally free $\mathcal O_S$-module $E$, the canonical
isomorphism $E\otimes E^\vee\simeq\mathcal{E}nd(E)$ sends
$\sum_i e_i\otimes e_i^\vee$ to $\operatorname{id}_E$.
A basis change and its inverse dual cancel in this tensor, so it
is independent of the local basis.  Apply the two universal quotient
maps to this identity tensor.  This defines $\delta$, with local
equation $\sum_iX_iY_i$, and shows directly that it commutes with
base change.  On a test scheme, let the two quotients be
$E\twoheadrightarrow L$ and $E^\vee\twoheadrightarrow M$.
The section corresponds to the composite
$M^\vee\hookrightarrow E\twoheadrightarrow L$.
Its vanishing is precisely incidence, proving the description of
the zero scheme.

On $P=\mathbf P(E)$ the kernel $K$ of the universal quotient
$\pi^*E\twoheadrightarrow\mathcal O_P(1)$ is locally free of
rank $m-1$.  The quotient splits locally.  On such an open, choose
a basis with the quotient equal to the last-coordinate projection.
The incidence locus in the dual projective bundle is then
$\mathbf P(K^\vee)$, cut out by the last homogeneous coordinate.
On an affine projective chart the equation is either a polynomial
coordinate, which is a nonzerodivisor over any coefficient ring,
or a unit, in which case its zero scheme is empty.  The displayed
two-term complex is therefore exact on every chart and hence globally.

For several dual factors, take product affine charts after the same
local splitting over $P$.  If one incidence equation is a unit,
its two-term complex is contractible and the intersection is empty.
Otherwise the equations are distinct polynomial coordinates, one
from each factor.  Successive quotients remain polynomial rings
over the same coefficient ring, so each next coordinate is a
nonzerodivisor.  They form a regular sequence.  Inductively tensoring
the two-term free resolutions consequently resolves their common
quotient.  Exactness is local and gives the asserted global Koszul
resolution.  For $m=1$ each section is a unit and the same empty-locus
argument applies.
\end{proof}

Two further exact certificates supplement these proofs:
\begin{itemize}
\item \path{code/C_determinant_scope_certificate.py} records
35 passing assertions for the explicit derived-tensor dimensions and
the 16-matrix trace example.
\item \path{code/C_trace_invariance_certificate.py} records
22 passing assertions, including all 16 simultaneous-centralizer cases.
\end{itemize}
The certificates do not replace the general proofs above.

These corrections do not assert a failure of the determinant
construction.  In using its extended axioms one must retain the
original's strict diagrams of actual complexes: the footnote to
Definition~4 explicitly rules out replacing that condition by
arbitrary commutativity in the derived category.  The later extensions
to ordinary triangles have the additional perfect-cohomology or
reduced-base hypotheses stated there
\cite[pp.~37--46, especially p.~41 footnote and pp.~43--44]{ROOT_KnudsenMumford1976}.

